% ===========================================================================
\subsection{\file{new-design-reference/assets/ASSETS.md}}
% ===========================================================================

\textbf{Markdown, 32 lines.} A checklist of the eleven files somebody would need to
produce before this design could be considered finished, plus three notes about
placeholder content to replace.

\begin{provenance}{new-design-reference/assets/ASSETS.md}
Hand-written, added in the first commit \code{1f7ed12} alongside the design it
describes, and never edited. Written by the same person in the same sitting as the
HTML, which is evident from its references to the exact class names used there
(\code{.hero-photo-placeholder}, \code{wt1} through \code{wt6}, \code{cv1},
\code{cv2}) and from the fact that line 9 quotes, character for character, the
\code{<img>} tag already written into the HTML's banner comment at line 34.
\textbf{Confidence: high, from the commit history and the internal references.}
\end{provenance}

\subsubsection*{Why it exists}

It is a handover note. The design was built with placeholder visuals, and this file
records precisely what would replace each one and what HTML change each replacement
needs. As an artefact it is genuinely good practice: it converts ``this is not
finished'' from a feeling into a list.

It is also the document that makes the state of this project unambiguous. The
placeholder in the hero could be argued to be a design choice; a checklist listing
the photograph as \emph{required before launch} cannot.

\begin{jargon}{Markdown}
A way of writing formatted text in a plain file, using punctuation instead of menus:
\code{\#} starts a heading, backticks mark code, and pipe characters draw a table. It
is designed to be readable as plain text and is the standard format for notes kept
alongside code.
\end{jargon}

\subsubsection*{Line by line}

\begin{lstlisting}[firstnumber=1,caption={\file{ASSETS.md}, lines 1--19}]
# Assets checklist: salmanadnan.com

Drop real assets here, then update the HTML as noted.

## Required before launch

| File | Size | Where used | HTML change needed |
|---|---|---|---|
| `salman.jpg` | 1200px tall min, 2:3 ratio | Hero photo left side | Replace `.hero-photo-placeholder` div with `<img src="/assets/salman.jpg" alt="Salman Adnan" style="width:100%;height:100%;object-fit:cover;object-position:top;">` |
| `thumbs/work-1.jpg` | 800x450px | Work grid card 1 | Add `style="background-image:url('/assets/thumbs/work-1.jpg')"` to `.wt1` |
| `thumbs/work-2.jpg` | 800x450px | Work grid card 2 | Same pattern for `.wt2` through `.wt6` |
| `thumbs/work-3.jpg` | 800x450px | Work grid card 3 | |
| `thumbs/work-4.jpg` | 800x450px | Work grid card 4 | |
| `thumbs/work-5.jpg` | 800x450px | Work grid card 5 | |
| `thumbs/work-6.jpg` | 800x450px | Work grid card 6 | |
| `case-1.jpg` | 1200x750px | Case study 1 visual | Add `background-image` to `.cv1` |
| `case-2.jpg` | 1200x750px | Case study 2 visual | Add `background-image` to `.cv2` |
| `favicon.svg` | 32x32 | Browser tab | Add `<link rel="icon" href="/assets/favicon.svg">` in `<head>` |
| `og-image.png` | 1200x630px | Social share | Update `og:image` meta tag |
\end{lstlisting}

Line 1 is the title and line 3 the instruction: put the real files in this folder,
then change the HTML as the table says.

Lines 7 and 8 open a table, line 8 being the separator row Markdown requires between
a table's headings and its body.

The eleven rows are the substance, and reading them tells you exactly what the
design is pretending to have.

\textbf{The photograph} is first, and it is the one that matters. Line 9 specifies
it precisely: at least 1200 pixels tall, in a 2:3 portrait ratio, and it gives the
complete replacement tag including \code{object-fit: cover} so the image fills the
frame without distorting, and \code{object-position: top} so a face stays in shot
when the frame crops. Whoever wrote this knew exactly what was needed and did not
have the file.

\textbf{Six project thumbnails} at 800 by 450, which is exactly the 16 by 9 shape the
stylesheet reserves with \code{aspect-ratio} at line 341. Line 11 gives the pattern
for the remaining five rather than repeating it, which is why rows 12 to 15 have an
empty final column.

\textbf{Two case-study images} at 1200 by 750, matching the \code{16/10} aspect ratio
set at \file{site.css} line 441. That the two ratios differ, and that each listed
asset matches its own, is a small sign of care.

\textbf{A favicon} and \textbf{a social preview image}, each with the exact tag
needed to install it.

\begin{honest}{None of the eleven exist, and the live site has at least one of them}
The \file{assets/} folder contains exactly one file: this checklist. Every asset it
calls ``required before launch'' is absent, so the design was never finished to the
standard its own author set.

The comparison that makes this sharp: fetching the live salmanadnan.com on
2026-08-11 found \code{<meta property="og:image" content=
"https://salmanadnan.com/assets/og.png">} along with explicit width and height
tags. So the social preview image listed here as required does exist, for the build
that replaced this one. This design was superseded by a version that finished the
checklist.
\end{honest}

\begin{lstlisting}[firstnumber=20,caption={\file{ASSETS.md}, lines 20--32}]

## Numbers to update

Search for `data-target` in `index.html` (metrics band) and replace all four values with real numbers.
Search for the three `hero-mini-num` divs in the hero and update those too.

## Booking link

Search for `YOUR_CAL_LINK` in `index.html` and replace with the actual booking URL.

## Client names

Three `<!-- UPDATE -->` comments in the work grid and case study sections mark placeholder client entries.
\end{lstlisting}

Line 23 is the metrics instruction, and line 24 is the one worth noticing.

\begin{keyidea}{This file knows about the duplication and tells you to fix it in
both places}
Line 23 says to search for \code{data-target} in the metrics band. Line 24 then says
to search for the three \code{hero-mini-num} divs and update those too.

That is the author documenting, at the time of writing, that the same three numbers
live in two places in two formats, as section 7.1 described. It was a recognised
cost, accepted deliberately, and written down. That is a materially different thing
from an accident, and it is why this document treats the duplication as a design
trade rather than a defect.

The sibling repository has no equivalent note because it has no hero statistics, and
consequently no fallback when JavaScript does not run.
\end{keyidea}

Line 28 tells the reader to search for \code{YOUR\_CAL\_LINK}.

\begin{honest}{Two of this file's three instructions are stale}
Searching \file{index.html} for \code{YOUR\_CAL\_LINK} finds nothing. The literal
placeholder was replaced on 2026-07-28 by commit \code{014406e}, which set the link
to \code{cal.com/salman-adnan/meeting}, a real personal calendar. So somebody
following this instruction searches, finds nothing, and cannot tell whether the job
is done or the instruction is wrong. Here the job \emph{is} done, unlike in the
sibling repository where the same stale instruction hides a link still pointing at
another project's calendar.

Line 32 is stale in a different way. It refers to ``three \code{<!-{}- UPDATE -{}->}
comments in the work grid and case study sections'' marking placeholder client
entries. There are no such comments anywhere in \file{index.html}: its
\code{UPDATE} notes live in the banner comments at the top of each section and none
of them is about client names. Either they were removed or they were never written.

Only the metrics instruction on lines 23 and 24 is verifiably still accurate, and
the \code{UPDATE} comment at \file{index.html} line 256 agrees with it.
\end{honest}

\subsubsection*{What it connects to}

It describes \file{new-design-reference/index.html} and
\file{new-design-reference/css/site.css} by name and by class, and it is the only
file in the repository documenting the project's intent. Nothing reads it: it is for
people.

The deploy workflow explicitly excludes it from publication (line 82 of the
workflow), which is the right instinct, and that exclusion no longer matches its
path. See section 7.6.

\subsubsection*{Concepts introduced}

Markdown tables; \code{object-fit} and \code{object-position}, the two properties
that control how an image fills a frame it does not exactly match; and the standard
sizes for a favicon (32 by 32) and an Open Graph preview image (1200 by 630).

\subsubsection*{How it breaks}

As a document it cannot break. As a source of truth it already has: two of its three
instructions describe a state of the code that no longer exists or never did, which
is the characteristic failure of documentation kept beside code and not updated with
it.

\subsubsection*{Honest findings for this file}

Two stale instructions (lines 28 and 32); eleven assets specified in useful detail
and never delivered; and, to its credit, one instruction (lines 23 and 24) that is
still exactly right and documents a known duplication rather than hiding it.
